\section{Methodology}\label{sec:method}

\subsection{The game}\label{sec:game}

One agent runs a stand for \NRounds{} rounds. Each round it chooses one of three locations, a sugar level, a
lemon level, one of two colours, and a price. Sugar, lemon and price run from 0.0 to 10.0 in steps of 0.1. The
agent then receives that round's profit and one line of customer feedback. Profit is deterministic. Each
location in a world has a hidden ideal setup and a hidden ceiling, and the profit for a setup is the
ceiling minus a linear penalty for every unit of distance from the ideal on each of sugar, lemon and price,
minus a flat penalty if the colour is wrong, floored at zero. The same setup in the same world always pays
the same profit. The feedback line picks one of sugar, lemon or price at random and says whether it is above,
below or at the ideal for the location the agent sold at. It is a local signal about one dial at one
location. It says nothing about other locations and never mentions colour. The agent never sees the profit
rule or the ideals; it has to find a good setup by trying things. Appendix~\ref{app:prompts} gives the
prompts verbatim.

\subsection{Worlds}\label{sec:worlds}

A world is a set of hidden ideals, one per location, together with each location's ceiling and penalty
slopes. Worlds are drawn by a generator that reproduces the structure of the original game
\citep{ederer2013pay}: one \emph{anchor} location with a ceiling between 90 and 110 and steep penalties, one
\emph{best} location with a ceiling between 180 and 220, equally steep penalties and a large colour penalty,
and a third location with a low ceiling and gentle penalties. A generated world is kept only if moving to the
best location without also changing the recipe pays less than staying at the anchor, so that reaching the top
requires a coordinated change of several variables and passes through a valley. The eight worlds used here
are the first eight of the validation split of the generator; no test world is used. Their global optima, the
most any setup can pay in one round, are listed in Appendix~\ref{app:worlds}. All scores below are expressed
as a share of the world's global optimum so that worlds can be compared. No setup outside the best location
pays more than 56.3\% of the optimum in any of the eight worlds.

\subsection{Two framings}\label{sec:framings}

\paragraph{Lemonade framing.}
The agent is told it runs a lemonade stand. Locations have their real names, such as school or stadium, and
the choices are called sugar, lemon, colour and price. The agent also receives a letter from the previous
manager describing the setup that worked for him at one location, with its profit, and saying that he never
tried another location and that other locations may need a very different recipe. The letter's setup is near,
but not at, the anchor location's ideal, and the anchor location is never the best one.

\paragraph{Alien framing.}
The worlds and their numbers are identical. Everything else is stripped away. The agent is told it runs a
business with five decision variables, A to E. Variable A takes the values A1, A2 and A3, which are the three
locations; B, C and E are sugar, lemon and price; D takes the values D1 and D2, the two colours. Profit is
called payoff. There is no letter and no recommended starting setup. Feedback has the same content as in the
lemonade framing, phrased as ``at this value of A, B is higher than ideal''. The alien framing was designed
following \citet{billinger2014search} and the label manipulation of \citet{austin2026semantic}: keep the
landscape, change only the words a language model has priors about.

\subsection{Teams and the shared board}\label{sec:board}

A team is \NAgents{} agents of the same model in the same world. Each agent runs its own stand in its own copy
of the world, so profits are not coupled and every number stays comparable with a single agent. Rounds are in
lockstep. When sharing is on, a shared board is appended to each agent's prompt listing, for every teammate and
every round so far, the setup chosen and the profit earned. In the default conditions all three agents choose
simultaneously from the board as it stood at the end of the previous round. Results are posted to the board
from round one in every condition, so a board that opens late already carries the private history. The
agents are not told how they are scored. The task instruction in every condition is ``Your goal is to
maximise the overall performance of the business'' (or ``of the lemonade stand''). Two details of the board's
wording are disclosed in Section~\ref{sec:limitations}.

\subsection{Conditions}\label{sec:conditions}

Nine conditions were run. They differ only in when the board is visible and, in one case, in what the agents
are told about it. Everything else is identical.

\begin{table}[H]
  \centering
  \caption{The nine conditions. Each was run in both framings on the same eight worlds.}
  \label{tab:conditions}
  \small
  \renewcommand{\arraystretch}{1.35}
  \begin{tabular}{p{4.6cm} p{10.2cm}}
    \toprule
    Condition & What happens \\
    \midrule
    Working alone & No shared board. Three agents play the same world without seeing each other. \\
    Sharing all along & The board is visible from round 1. \\
    Sharing after 2 private rounds & The board opens in round 3. \\
    Sharing after 5 private rounds & The board opens in round 6. \\
    Sharing after 10 private rounds & The board opens in round 11. \\
    Sharing after 15 private rounds & The board opens in round 16. \\
    Sharing at rounds 6, 11 and 16 only & The board is visible only in those three rounds, five rounds apart, and private otherwise. We call this the periodic board. \\
    Sharing, told to stay different & The board is visible from round 1. Each agent is also told to pick an approach different from the ones already claimed, to adopt a teammate's approach only when its posted result is clearly better, and to keep one variation of its own when it does. \\
    One agent at a time & The board is visible from round 1. Within each round the agents move in a rotating order (1-2-3, then 2-3-1, then 3-1-2), so the second mover sees the first mover's result from the current round and the third sees both. \\
    \bottomrule
  \end{tabular}
\end{table}

Working alone is the reference. Three agents working alone receive the same \NAttempts{} attempts in a world
as a sharing team, so the two are scored on the same number of tries. Where the text says the periodic board it
means the condition in which the board is visible only at rounds 6, 11 and 16.

\subsection{Measures}\label{sec:measures}

All measures are computed per team per world and expressed as a share of the world's global optimum, except
the count of distinct setups.

\begin{itemize}[leftmargin=*, itemsep=2pt]
  \item \textbf{Maximum single-round profit (max):} the highest profit any one agent on the team reached in
    any round, as a share of the global optimum. For working alone this is the best of the three independent
    agents, so both arms are scored on the best of \NAttempts{} attempts. This measures whether the team
    found the best setup.
  \item \textbf{Cumulative earnings:} everything the three agents earned over all \NRounds{} rounds, as a
    share of what they would have earned by hitting the best setup in every round. This measures what the
    team earned while it was still searching.
  \item \textbf{Distinct setups:} the number of different full setups (location, sugar, lemon, colour, price)
    among the team's \NAttempts{} attempts, pooled across the three agents. This measures how much the team
    explored.
  \item \textbf{Reached the best location:} whether the team's max exceeds 60\% of the optimum, which only
    a setup at the best location can do.
\end{itemize}

The first two measures are the two goals a team might have, best-arm identification and regret minimisation,
and the agents are instructed on neither.

\subsection{Model, runs and statistics}\label{sec:stats}

The main runs use GPT-5.6 (gpt-5.6-sol) through the OpenAI API with the provider's default sampling. Each
condition was run on the eight worlds in both framings. Three follow-ups, each planned before it was run, are
reported in Section~\ref{sec:results}: working alone and the periodic board on sixteen further worlds in both
framings, three conditions on a second model, Claude Opus 5.5, and working alone and sharing all along in
60-round games on the eight worlds. The first run of every condition was made in October 2026. To
test whether the first run's numbers were a fluke, every condition was then repeated under a plan written
before any repeat was run: working alone, sharing all along and sharing at rounds 6, 11 and 16 were repeated
twice more, giving three runs per world, and the other six conditions were repeated once, giving two runs per
world, in both framings. Repeats use their own environment seeds for the customer-feedback draws. Every planned
episode completed, and no run was stopped, extended or excluded on the basis of its result.

The unit of analysis is the world. A world's value for a condition is the mean over its runs. Condition means
in the tables are the mean over the eight worlds of those values, and the figure after $\pm$ is the sample
standard deviation of the condition's run means across runs, which shows the run-to-run wobble. For each
contrast we take the paired difference between a sharing condition and working alone within each world, and
report the mean, a 95\% $t$ interval over the eight world-level differences, and the exact two-sided sign-flip
permutation test over the eight worlds. The smallest $p$ this test can produce with eight worlds is
$0.0078$. The headline contrast is the mean of the eight sharing conditions minus working alone, reported
for the first run alone, for the repeat runs alone and for all runs, so that the reader can see whether the
first run and the repeats agree.
